\section{Fase 3: arquitecturas e implementación}
\label{sec:fase3}

A la fecha están implementadas y evaluadas las cuatro arquitecturas: la de
referencia (\eegnet{}), el Conformer, el modelo fundacional LaBraM (variante Base) y
EEGMamba, cuya comparación se presenta en la sección~\ref{sec:comparacion}. EEGMamba
requirió las adaptaciones específicas descritas en la
sección~\ref{sec:eegmamba-adapt}.

\subsection{Representación de entrada y procesamiento por arquitectura}
\label{sec:representaciones}

Las cuatro arquitecturas reciben la misma señal preprocesada, pero no la consumen de la
misma forma. Esta sección detalla qué hace cada una con ella, porque la
interpretación de los resultados depende de ello: dos de las cuatro operan sobre la forma
de onda continua y las otras dos la fragmentan en parches de un segundo. Las dimensiones que
se indican, $62$ canales por $801$ muestras, son las de los protocolos sobre \openbmi{}; en
el capítulo~\ref{cap:validacion-externa} la entrada tiene los $45$ canales comunes a los
dos conjuntos, y el número de parámetros de la capa que depende de los canales cambia en
consecuencia.

\subsubsection{\eegnet{}: convolución separable sobre la señal continua}

\eegnet{} \cite{lawhern2018}, en la implementación de \texttt{braindecode} 1.6.1
\cite{schirrmeister2017} con sus hiperparámetros por defecto, consume el tensor $(C \times T) = (62 \times 801)$ sin ninguna transformación intermedia.
Su primer bloque aplica convoluciones \emph{temporales}, cuyos núcleos actúan como filtros
de banda aprendidos; el segundo aplica una convolución \emph{en profundidad} sobre el eje
de canales, un filtro espacial por cada filtro temporal, que cumple un papel análogo al de
los patrones espaciales comunes; el tercero combina ambos ejes con una convolución
separable. Un promediado temporal y una capa lineal producen los logits
(figura~\ref{fig:arq-eegnet}). Con $62$ canales
tiene $2\,898$ parámetros, de dos a algo más de tres órdenes de magnitud menos que las demás.

\begin{figure}[H]
\centering
\begin{tikzpicture}[
  font=\scriptsize,
  b/.style={draw, rounded corners=2pt, minimum height=0.7cm, align=center, fill=blue!8},
  fl/.style={-{Latex[length=1.6mm]}, semithick}, x=1cm, y=1cm
]
\node[b, minimum width=1.9cm] (in) at (0,0) {$62 \times 801$\\señal continua};
\node[b, minimum width=2.0cm] (t)  at (2.6,0) {conv.\ temporal\\(filtros de banda)};
\node[b, minimum width=2.0cm] (e)  at (5.3,0) {conv.\ en profundidad\\(filtros espaciales)};
\node[b, minimum width=1.9cm] (sp) at (7.9,0) {conv.\ separable\\+ promediado};
\node[b, minimum width=1.2cm, fill=orange!18] (o) at (10.0,0) {lineal\\2 clases};
\foreach \a/\bb in {in/t,t/e,e/sp,sp/o} \draw[fl] (\a) -- (\bb);
\end{tikzpicture}
\caption{\eegnet{}. Opera sobre la señal continua; separa el filtrado temporal del
espacial, en ese orden.}
\label{fig:arq-eegnet}
\end{figure}

\subsubsection{Conformer: convolución seguida de atención}

El EEG-Conformer \cite{song2023conformer}, también en la implementación de
\texttt{braindecode} con sus hiperparámetros por defecto, recibe el mismo tensor
$(62 \times 801)$ que \eegnet{} y no necesita adaptador. Un bloque
convolucional inicial, temporal y después espacial, transforma la señal en una secuencia de
vectores o \emph{tokens}; sobre esa secuencia actúa un codificador con autoatención, que
pondera cada instante en función de todos los demás en lugar de hacerlo dentro de una
ventana local fija como la convolución. Un perceptrón produce los logits. Tiene $710\,146$
parámetros (figura~\ref{fig:arq-conformer}).

\begin{figure}[H]
\centering
\begin{tikzpicture}[
  font=\scriptsize,
  b/.style={draw, rounded corners=2pt, minimum height=0.7cm, align=center, fill=blue!8},
  fl/.style={-{Latex[length=1.6mm]}, semithick}, x=1cm, y=1cm
]
\node[b, minimum width=1.9cm] (in) at (0,0) {$62 \times 801$\\señal continua};
\node[b, minimum width=2.4cm] (c)  at (2.9,0) {conv.\ temporal\\+ conv.\ espacial};
\node[b, minimum width=1.9cm, fill=green!12] (tk) at (5.7,0) {secuencia de\\\emph{tokens}};
\node[b, minimum width=2.1cm] (at) at (8.2,0) {autoatención\\(global)};
\node[b, minimum width=1.2cm, fill=orange!18] (o) at (10.4,0) {MLP\\2 clases};
\foreach \a/\bb in {in/c,c/tk,tk/at,at/o} \draw[fl] (\a) -- (\bb);
\end{tikzpicture}
\caption{Conformer. La convolución construye los \emph{tokens} y la autoatención los
relaciona entre sí sin restricción de distancia temporal.}
\label{fig:arq-conformer}
\end{figure}

\subsubsection{LaBraM: parches de un segundo con identidad de canal}

LaBraM \cite{jiang2024labram} es el único modelo con pesos preentrenados. Divide la ventana
en parches de $200$ muestras,
es decir $1$\,s a \hz{200}, lo que obliga a truncar de $801$ a $800$ y produce $4$ parches
por canal. Cada parche se proyecta a un vector, al que se suma un \textbf{embedding de la
identidad del canal}: la entrada incluye explícitamente el electrodo de origen
de cada parche, información que no forma parte de la entrada de las tres arquitecturas restantes. Esa lista
de índices se toma de un vocabulario de $128$ posiciones de electrodo, que incluye
posiciones de los sistemas 10-10 y 10-05, y es lo que permite aplicar el
modelo a montajes distintos. Un codificador \emph{transformer} procesa el conjunto de
parches. Se emplea la variante Base, $5\,818\,138$ parámetros, con ajuste fino completo,
partiendo de los pesos distribuidos en \texttt{braindecode/labram-pretrained}
(figura~\ref{fig:arq-labram}).

Tres decisiones de esta adaptación se apartan del procedimiento con que se preentrenó el
modelo y condicionan su evaluación. La señal se le entrega normalizada por canal
(\textit{z-score}), igual que a las demás arquitecturas, y no en microvoltios divididos por
$100$ como en el preentrenamiento. El \emph{embedding} temporal preentrenado, de $16$
posiciones, se interpola a las $5$ que exige la ventana de $4$\,s. La capa de clasificación
se reinicia, y el ajuste fino usa una tasa de aprendizaje uniforme de $10^{-4}$, sin el
decaimiento por capa del procedimiento original.

\begin{figure}[H]
\centering
\begin{tikzpicture}[
  font=\scriptsize,
  b/.style={draw, rounded corners=2pt, minimum height=0.7cm, align=center, fill=blue!8},
  fl/.style={-{Latex[length=1.6mm]}, semithick}, x=1cm, y=1cm
]
\node[b, minimum width=1.9cm] (in) at (0,0) {$62 \times 801$\\señal continua};
\node[b, minimum width=2.2cm, fill=green!12] (pa) at (2.9,0)
  {parches de \seg{1}\\$62 \times 4 \times 200$};
\node[b, minimum width=2.0cm] (em) at (5.7,0) {proyección\\de cada parche};
\node[b, minimum width=2.0cm] (tr) at (8.3,0) {\emph{transformer}\\preentrenado};
\node[b, minimum width=1.2cm, fill=orange!18] (o) at (10.5,0) {lineal\\2 clases};
\node[b, minimum width=2.4cm, fill=red!10] (ch) at (5.7,-1.4)
  {índices de canal\\(vocabulario de 128)};
\foreach \a/\bb in {in/pa,pa/em,em/tr,tr/o} \draw[fl] (\a) -- (\bb);
\draw[fl] (ch) -- (em);
\end{tikzpicture}
\caption{LaBraM. Único modelo cuya entrada incluye la identidad del electrodo junto a la señal, lo
que le permite operar sobre montajes distintos del de preentrenamiento.}
\label{fig:arq-labram}
\end{figure}

\subsubsection{EEGMamba: espacio de estados sobre parches, con rama espectral}

Se emplea la implementación pública de EEGMamba \cite{wang2025eegmamba}, sin los pesos del
modelo fundacional que describe ese trabajo. Fragmenta la ventana en los mismos $4$
parches de $200$ muestras. Cada parche se codifica
por dos vías que se suman: una \emph{temporal}, una convolución de $25$ filtros de
longitud $49$ y paso $25$, y una
\emph{espectral}, que toma el módulo de la transformada rápida de Fourier real del parche,
$101$ coeficientes, y lo proyecta al mismo espacio; se añade una codificación posicional
convolucional. Es la única de las cuatro cuya entrada incluye
una representación en frecuencia de forma explícita. La secuencia resultante se procesa
mediante una pila de $12$ bloques Mamba2 \cite{gu2023mamba,dao2024mamba2}, modelos de
espacio de estados aplicados en ambos sentidos de la secuencia, cuyo
costo crece de forma lineal con la longitud de la secuencia y no cuadrática como el de la
autoatención. Se emplea la cabeza que conserva la representación por parche y canal;
tiene $3\,376\,445$ parámetros y se entrena desde cero
(figura~\ref{fig:arq-eegmamba}).

\begin{figure}[H]
\centering
\begin{tikzpicture}[
  font=\scriptsize,
  b/.style={draw, rounded corners=2pt, minimum height=0.7cm, align=center, fill=blue!8},
  fl/.style={-{Latex[length=1.6mm]}, semithick}, x=1cm, y=1cm
]
\node[b, minimum width=1.7cm] (in) at (0,0) {$62 \times 801$\\señal continua};
\node[b, minimum width=2.0cm, fill=green!12] (pa) at (2.6,0)
  {parches de \seg{1}\\$62 \times 4 \times 200$};
\node[b, minimum width=1.7cm] (tt) at (5.1,0.75) {proyección\\temporal};
\node[b, minimum width=1.7cm, fill=red!10] (ff) at (5.1,-0.75) {$|\text{rFFT}|$\\(101 coef.)};
\node[b, minimum width=2.2cm] (mm) at (7.9,0) {Mamba2\\bidireccional};
\node[b, minimum width=1.2cm, fill=orange!18] (o) at (10.2,0) {lineal\\2 clases};
\draw[fl] (in) -- (pa);
\draw[fl] (pa) |- (tt);
\draw[fl] (pa) |- (ff);
\draw[fl] (tt) -| (mm);
\draw[fl] (ff) -| (mm);
\draw[fl] (mm) -- (o);
\node[font=\scriptsize] at (6.6,0) {$+$};
\end{tikzpicture}
\caption{EEGMamba. Suma una codificación temporal y una espectral de cada parche, y procesa
la secuencia con bloques de espacio de estados en ambos sentidos.}
\label{fig:arq-eegmamba}
\end{figure}

Resumiendo el contraste que importa para la comparación: \eegnet{} y el Conformer operan
sobre la señal continua, mientras que LaBraM y EEGMamba la fragmentan en parches de un
segundo. LaBraM es el único cuya entrada incluye la identidad del electrodo y el único preentrenado;
EEGMamba es el único cuya entrada incluye una representación en frecuencia explícita. Ninguna de las
cuatro abandona el dominio del tiempo como espacio de predicción, punto sobre el que vuelve
la discusión.

\subsection{Esquema de entrenamiento unificado}

Todas las arquitecturas comparten el mismo esquema, con una excepción que se declara
más abajo (el recorte de la norma del gradiente), para que la comparación
atribuya las diferencias al modelo y no al ajuste: optimizador AdamW
\cite{loshchilov2019adamw} con decaimiento de pesos de $0.01$,
calentamiento lineal de $10$ épocas seguido de decaimiento coseno calculado sobre el
máximo de $200$ épocas, suavizado de etiquetas de $0.1$ en la pérdida de entrenamiento
(la de validación se calcula sin suavizado),
precisión mixta automática, tamaño de lote efectivo fijo \emph{dentro de cada
arquitectura}, y parada temprana por pérdida de validación con paciencia de 20 épocas
sobre un máximo de 200. Al terminar se restauran los pesos de la época de menor pérdida
de validación.

Dos consecuencias de esta configuración no son evidentes y afectan a la lectura de los
resultados. La paciencia no corre durante el calentamiento, de modo que la parada más
temprana posible es la época $30$. Y como el coseno se calcula sobre $200$ épocas mientras
las ejecuciones \loso{} se detienen entre las épocas $30$ y $146$, la tasa de aprendizaje
decae poco antes de la parada en la mayoría de ellas: conserva el $97$\,\% de su valor en la
época $30$ y el $19$\,\% en la $146$. El cómputo no es determinista bit a bit (cuDNN sin modo determinista y precisión
mixta), así que repetir una ejecución con la misma semilla puede variar en los últimos
decimales. Las versiones empleadas son \texttt{braindecode} 1.6.1, MNE-Python 1.12.1,
\texttt{mne-icalabel} 0.9.0 y MOABB 1.5.0 \cite{jayaram2018moabb}. Las ejecuciones se
agrupan en dos regímenes de entrenamiento. El \emph{primer régimen} reúne las ejecuciones
previas (GPU NVIDIA A40; recorte de la norma del gradiente solo en EEGMamba); el
\emph{segundo régimen}, las posteriores (GPU NVIDIA L40S; recorte a $1.0$ en las cuatro
arquitecturas). La
tabla~\ref{tab:regimen} indica el régimen de cada semilla y la tabla~\ref{tab:hardware} el
hardware; esta denominación se usa en el resto del documento.

\textbf{Recorte de la norma del gradiente.}\label{sec:recorte-gradiente}
En el primer régimen, las ejecuciones base intra-sujeto y \loso{} de \eegnet{}, el
Conformer y LaBraM se entrenaron \emph{sin} recorte de la norma del gradiente, y las de
EEGMamba con recorte a $1.0$. En el segundo régimen las cuatro arquitecturas usan recorte a
$1.0$. La transferencia de la semilla $2026$ también lo llevó, aunque se entrenó en A40
(tabla~\ref{tab:regimen}). Cada ejecución guarda un identificador de configuración
(\textit{hash}) que permite comprobar si el recorte estuvo activo.

El recorte se usa en el segundo régimen porque los núcleos de espacio de estados de EEGMamba
saturan el rango de \texttt{float16} y divergen sin él, y es una técnica descrita para los
gradientes explosivos \cite{pascanu2013difficulty}; aplicarlo a las cuatro arquitecturas
mantiene un único procedimiento de entrenamiento. No se comprobó que fuese neutro para
\eegnet{}, el Conformer y LaBraM, y no se conserva registro numérico de la divergencia de
EEGMamba sin recorte ni de la norma del gradiente, por lo que se desconoce en qué fracción de
los pasos el recorte actúa. Las recetas de referencia de las cuatro arquitecturas no
coinciden en usarlo: el repositorio de EEGMamba \cite{wang2025eegmamba} lo activa por defecto
con valor $1$; el de LaBraM \cite{jiang2024labram} lo deja desactivado en el ajuste fino; y
en los de \eegnet{} \cite{lawhern2018} y del Conformer \cite{song2023conformer} no aparece.
Este trabajo no afirma que el recorte sea una práctica habitual en el decodificado de EEG.

En \eegnet{}, el Conformer y LaBraM, el contraste entre una variante (segundo régimen) y la
ejecución base de la semilla $2026$ mezcla el efecto de la variante con un cambio compuesto
del régimen de entrenamiento: recorte de gradiente, hardware (A40 frente a L40S) y versión
del código. En EEGMamba el recorte no cambia, aunque sí el hardware y el código. Ese cambio
no puede aislarse en sus componentes, pero se acota su magnitud comparando, para cada sujeto,
la semilla $2026$ de la base con la media de las semillas $2027$ y $2028$
(sección~\ref{sec:cota-confusor}; resultados en la sección~\ref{sec:confusor}). No se detecta
efecto del cambio de régimen en ninguna arquitectura: el límite del intervalo de confianza al
$95$\,\% de la diferencia media por sujeto en exactitud es de $0.010$ a $0.020$ según la
arquitectura, y en una media de tres semillas, de las cuales una se entrenó sin recorte, la
contaminación máxima es de $0.003$ a $0.007$ (tabla~\ref{tab:confusor}). Eso descarta un
efecto medio mayor que esos límites, no un efecto nulo. Para el protocolo intra-sujeto de
\eegnet{} y del Conformer no existe base con recorte (las tres semillas se entrenaron sin él),
de modo que ese efecto no es estimable y se declara en las limitaciones. Repetir esa base con
recorte queda como trabajo pendiente.

\textbf{Alcance de la unificación.} Las cuatro arquitecturas comparten la \emph{estructura}
del procedimiento (optimizador, planificación de la tasa de aprendizaje, criterio de parada,
aumentación, particiones y métricas) y no los valores numéricos, que enumera la
tabla~\ref{tab:configs}.

\begin{table}[H]
\centering
\caption{Configuración efectiva de cada arquitectura. Las dos primeras columnas son los
valores con los que se ejecutaron los resultados reportados; la última resume por qué
difieren del resto. La tabla no recoge el recorte de la norma del gradiente, que no es un
valor por arquitectura sino un cambio de régimen entre ejecuciones (párrafo
``Recorte de la norma del gradiente'' de esta sección).}
\label{tab:configs}
\small
\begin{tabular}{lccp{7.2cm}}
\toprule
 & \textbf{Tasa de aprendizaje} & \textbf{Lote} & \textbf{Motivo de la desviación} \\
\midrule
\eegnet{}   & $2.8\times10^{-3}$ & $256$ & Empleada al pasar al lote $256$, sin búsqueda ni criterio de ajuste registrado; coincide con escalar $10^{-3}$ por la raíz cuadrada del aumento de lote, $10^{-3}\sqrt{8}$ \\
Conformer   & $1.0\times10^{-3}$ & $256$ & Valor por defecto del esquema; no se comparó de forma sistemática con el de \eegnet{} \\
LaBraM      & $1.0\times10^{-4}$ & $32$  & Parte de pesos preentrenados: en una prueba preliminar, con la tasa por defecto producía una salida constante para todas las entradas \\
EEGMamba    & $1.0\times10^{-3}$ & $32$  & Cómputo en \texttt{bfloat16}: los núcleos de espacio de estados superan el rango numérico representable en \texttt{float16} y divergen (sección~\ref{sec:eegmamba-adapt}) \\
\bottomrule
\end{tabular}
\end{table}

Los ajustes responden a problemas de entrenamiento observados en pruebas preliminares y no
a la búsqueda de un mejor resultado para una arquitectura concreta; no se realizó una
búsqueda sistemática de hiperparámetros. Sí hubo decisiones tomadas observando el conjunto
principal: la cabeza de EEGMamba se eligió comparando su exactitud intra-sujeto sobre las
mismas particiones que luego se reportan (sección~\ref{sec:eegmamba-adapt}) y se trasladó al
protocolo \loso{}, lo que pudo favorecer ligeramente a esa arquitectura, y la tasa de LaBraM
y la precisión \texttt{bfloat16} se fijaron tras pruebas preliminares sobre el mismo
conjunto. La comparación debe leerse, por tanto, entre arquitecturas en una configuración
funcional y no óptima. El tamaño de lote tampoco es constante ($256$ frente a $32$), lo que
altera el número de actualizaciones del optimizador por época y no solo la velocidad de
cómputo.

\textbf{Aumentación de datos.} Todos los entrenamientos, en las cuatro arquitecturas y
en los dos protocolos, aplican tres transformaciones sobre el conjunto de entrenamiento,
nunca sobre validación ni sobre prueba. Se aplican \emph{por minilote}, de modo que cada
época ve realizaciones distintas:

\begin{itemize}
  \item \textbf{Desplazamiento temporal} de hasta el $5$\,\% de la ventana, es decir
  $\pm 40$ muestras o $\pm 200$\,ms, sorteado de forma independiente por segmento.
  \item \textbf{Ruido gaussiano} de desviación igual al $2$\,\% de la desviación típica
  \emph{de cada canal}, calculada sobre el propio lote.
  \item \textbf{Enmascarado en frecuencia} de una banda aleatoria de hasta $2$\,Hz,
  aplicada sobre la transformada de Fourier y deshecha con la transformada inversa.
\end{itemize}

La mezcla de segmentos (\textit{mixup}) está implementada pero desactivada. La
aumentación es idéntica en las cuatro arquitecturas y es la línea base frente a la que se
contrastan las variantes de aumentación. Es \textbf{en línea}: el conjunto de entrenamiento
conserva su tamaño (por ejemplo, los $18\,360$ segmentos de cada \textit{fold} \loso{}) y
solo varía entre épocas la perturbación con que cada segmento se presenta a la red, de modo
que actúa como regularización y no amplía el número de ejemplos. Solo una de las tres
transformaciones añade ruido; el desplazamiento temporal es geométrico y el enmascarado en
frecuencia suprime una banda.

La arquitectura de referencia es \eegnet{} \cite{lawhern2018}, el punto de comparación
habitual en la literatura de imaginación motora.

\subsection{Adaptaciones específicas de EEGMamba}
\label{sec:eegmamba-adapt}

A diferencia de LaBraM, EEGMamba se entrena desde cero, sin pesos preentrenados.
Por ello se aparta del esquema unificado en tres puntos, todos necesarios para que
el modelo aprenda y ninguno orientado a favorecerlo frente al resto de
arquitecturas:

\begin{itemize}
  \item \textbf{Precisión de cómputo.} Los núcleos de espacio de estados del
  bloque Mamba (recorrido selectivo y exponenciales) superan el rango numérico
  representable en media precisión \texttt{float16} y producen valores no numéricos pasadas las primeras
  épocas. Se sustituye \texttt{float16} por \texttt{bfloat16}, que conserva el
  rango de exponente de la precisión simple y elimina el desbordamiento sin
  penalización de velocidad apreciable.
  \item \textbf{Recorte de gradiente.} La norma del gradiente se acota a $1.0$
  \cite{pascanu2013difficulty}. En las ejecuciones base del primer régimen es una
  adaptación propia de esta arquitectura; en el segundo régimen lo llevan las cuatro. Su
  justificación y su consecuencia para la comparación se tratan en el párrafo ``Recorte de
  la norma del gradiente'' de la subsección anterior.
  \item \textbf{Cabeza de clasificación.} Se emplea en \textbf{ambos protocolos} la
  cabeza que conserva la representación de todos los canales y parches y la proyecta
  linealmente, porque es la que preserva la lateralización sensoriomotora, el contraste
  que separa la imaginación motora de una mano y de la otra. La alternativa considerada,
  un promediado global sobre canales y parches, elimina la información de lateralización antes de la
  proyección, y con ella la exactitud intra-sujeto queda en $0.519$, apenas por encima
  del azar (Wilcoxon de una muestra frente a $0.50$, $p = 0.017$). Conservando la
  representación por canal asciende a $0.577$ con la semilla 2026 ($p = 1.5\times10^{-8}$ frente al azar) y
  supera a la variante de promediado con tamaño de efecto grande
  ($p = 2.6\times10^{-6}$, $r = 0.64$; contraste de Wilcoxon pareado sobre los 54
  sujetos).

  Esta cabeza sobreajusta en el protocolo intra-sujeto. Con 62 canales la proyección
  recibe $62 \times 4 \times 200 = 49\,600$ características frente a los 144 segmentos de
  entrenamiento de cada partición, y los registros por época muestran que la pérdida de
  validación aumenta desde las primeras épocas mientras la de entrenamiento disminuye.
  La consecuencia es que, en el $85.6$\,\% de las ejecuciones intra-sujeto, el criterio de
  parada conserva un modelo de las épocas de calentamiento
  (sección~\ref{sec:seleccion-intra}). La exactitud de $0.577$ debe leerse, por tanto,
  como el resultado conjunto de la arquitectura y del procedimiento de selección del
  modelo con 36 ensayos de validación, y no como una medida de capacidad comparable con
  la de las otras tres arquitecturas.
\end{itemize}

La tasa de aprendizaje se fija en $10^{-3}$, superior a la empleada al ajustar el
modelo preentrenado, coherente con un entrenamiento desde cero.

\subsection{Semillas y régimen de optimización de cada una}
\label{sec:semillas-metodo}\label{sec:fase3-variantes}

Las cuatro arquitecturas se ejecutaron con tres semillas, $2026$, $2027$ y $2028$, en el
protocolo intra-sujeto y en \loso{}, y las tres semillas se repitieron también en el modelo
de $54$ sujetos transferido a WBCIC-SHU, en el modelo nativo de WBCIC-SHU y en el protocolo
de sujetos apartados (44 más 10). Eso significa, por protocolo:

\begin{itemize}
  \item \textbf{Intra-sujeto.} \eegnet{} y el Conformer ya contaban con las tres semillas;
  para LaBraM y EEGMamba se añadieron la $2027$ y la $2028$. Cada arquitectura aporta
  $54 \times 20 \times 3 = 3\,240$ mediciones.
  \item \textbf{\loso{}.} La semilla $2026$ procede de las ejecuciones del primer régimen, y se
  añadieron la $2027$ y la $2028$ para las cuatro arquitecturas: $54 \times 3 = 162$
  \textit{folds} por arquitectura.
  \item \textbf{Transferencia y sujetos apartados.} Se añadieron las semillas $2027$ y
  $2028$ a las tres configuraciones que antes tenían la $2026$.
\end{itemize}

El número de semillas se limita a tres por el presupuesto de cómputo. Se reportan todas; no
se elige la de mejor desempeño, porque seleccionarla sobre el propio conjunto de prueba
sería un sesgo de selección. Las variantes de preprocesamiento y de aumentación
(secciones~\ref{sec:variantes-preproc} y~\ref{sec:variantes-aug}) usan solo la semilla
$2026$, para cubrir más variantes con una semilla cada una; por ello el desplazamiento medio
entre semillas no se incorpora en esos contrastes (sección~\ref{sec:ablaciones-f4}).

Las semillas controlan la inicialización de la red y los sorteos del entrenamiento. Se
emparejan por índice entre arquitecturas, aunque las inicializaciones son independientes, de
modo que el emparejamiento solo etiqueta. La semilla del remuestreo estadístico es otra
(sección~\ref{sec:bootstrap}).

\textbf{Régimen de optimización de cada semilla.} La tabla~\ref{tab:regimen} resume el
régimen de recorte de cada semilla, comprobado con el identificador de configuración de cada
ejecución, y la tabla~\ref{tab:hardware} el hardware en que se entrenó cada parte.

\begin{table}[H]
\centering
\caption{Régimen de recorte de gradiente por semilla (orden $2026$, $2027$, $2028$).
S: sin recorte. C: con recorte a $1.0$. En las semillas $2026$ de las cuatro arquitecturas,
las ejecuciones son del primer régimen (GPU A40); las demás, del segundo (GPU
L40S), salvo las semillas $2027$ y $2028$ intra-sujeto de \eegnet{} y del Conformer, que
también se entrenaron en A40. La transferencia de la semilla $2026$, aunque anterior al segundo régimen y
en A40, ya llevaba recorte.}
\label{tab:regimen}
\small
\begin{tabular}{lcccc}
\toprule
 & \textbf{\eegnet{}} & \textbf{Conformer} & \textbf{LaBraM} & \textbf{EEGMamba} \\
\midrule
\loso{}                          & S C C & S C C & S C C & C C C \\
Intra-sujeto                     & S S S & S S S & S C C & C C C \\
Transferencia y 44 + 10          & C C C & C C C & C C C & C C C \\
\bottomrule
\end{tabular}
\end{table}

La tabla muestra qué comparaciones cruzan regímenes: las de \eegnet{}, el Conformer y
LaBraM entre la semilla $2026$ y las otras dos en \loso{}, y las de LaBraM en
intra-sujeto. Entre las semillas de una misma fila sin cambio de letra no hay cambio de
recorte, aunque sí puede haber cambio de hardware.

\begin{table}[H]
\centering
\caption{Hardware de entrenamiento. En la A40, memoria, capacidad de cómputo y programas son
los que reportan los registros de las ejecuciones; en la L40S, y para la arquitectura de
ambas GPU, son los de catálogo y los de la imagen del contenedor empleada. No
se interpreta ninguna diferencia entre ejecuciones como efecto del hardware, y los tiempos por
\textit{fold} no son comparables entre las dos GPU.}
\label{tab:hardware}
\small
\begin{tabular}{@{}p{3.6cm}p{5.2cm}p{5.2cm}@{}}
\toprule
 & \textbf{NVIDIA A40} & \textbf{NVIDIA L40S} \\
\midrule
Arquitectura (capacidad de cómputo) & Ampere (sm\_86) & Ada Lovelace (sm\_89) \\
Memoria de la GPU                   & $48$\,GB ($47.7$\,GB útiles) & $48$\,GB \\
Régimen de entrenamiento            & Primero & Segundo \\
Qué se entrenó                      & Ejecuciones base de la semilla $2026$ de las cuatro arquitecturas (intra-sujeto y \loso{}), semillas $2027$ y $2028$ intra-sujeto de \eegnet{} y del Conformer, y transferencia de la semilla $2026$ & Las demás semillas y todas las variantes (ablaciones, aumentación, alineamiento, calibración) \\
Programas                           & PyTorch $2.8.0$, CUDA $12.8$ & PyTorch $2.8.0$, CUDA $12.8$ \\
\bottomrule
\end{tabular}
\end{table}

\subsection{Variantes de aumentación de datos}
\label{sec:variantes-aug}

La configuración base aplica las tres transformaciones descritas en la
sección~\ref{sec:fase3} (desplazamiento temporal, ruido gaussiano y enmascarado en
frecuencia). Se evalúan tres cambios, todos sobre el conjunto de entrenamiento y por
minilote, como la base:

\begin{itemize}
  \item \textbf{Sin aumentación.} Las tres transformaciones de la base en cero. Esta
  condición no es la base: es la ausencia de toda aumentación.
  \item \textbf{Reflexión de canales} \cite{wang2024channelreflection}. La base más una
  reflexión que, con probabilidad $0.5$ por segmento, intercambia cada canal con su simétrico
  del otro hemisferio y permuta la etiqueta (mano izquierda por mano derecha), porque la
  desincronización de la imaginación motora es contralateral. El simétrico de un canal impar
  $n$ es el par $n+1$, según el sistema 10-20 y sus extensiones; los canales de la línea media
  y los que no tienen pareja en el montaje se quedan en su sitio. Se aplica antes que el
  resto de transformaciones del lote.
  \item \textbf{Ruido en la amplitud de Fourier.} La base más una perturbación
  multiplicativa gaussiana del módulo de cada componente de la transformada de Fourier,
  $|F'(f)| = |F(f)|\,\max\!\bigl(0,\ 1 + \sigma\,\mathcal{N}(0,1)\bigr)$ con $\sigma = 0.1$,
  independiente por canal, segmento y componente, con la fase conservada, y se vuelve al
  tiempo con la transformada inversa. No se dispuso de un valor de referencia en la
  literatura para esta aplicación; $\sigma = 0.1$ es una elección propia, fijada antes de
  ejecutar y sin ajuste posterior, por lo que el resultado vale para esa amplitud y no para
  el ruido en Fourier en general.
\end{itemize}

En el protocolo intra-sujeto se ejecutan las tres con la semilla $2026$; en \loso{}, solo la
reflexión de canales, con la semilla $2026$. El ruido en Fourier en \loso{} no se ejecutó.

\subsection{Variantes con una sola semilla y ablaciones recortadas en LaBraM y EEGMamba}
\label{sec:ablaciones-recortadas}

LaBraM y EEGMamba cuestan varias horas cada ablación intra-sujeto, porque
saturan la GPU y la paralelización por trabajadores apenas las acelera. Por coste de
cómputo se decidió, antes de ejecutarlas, limitar sus ablaciones a la reflexión de canales,
el ruido en Fourier y la ventana de \seg{0} a \seg{2}, esta última porque cambia el número
de parches. \eegnet{} y el Conformer recorren las ocho: tres subbandas, dos ventanas, sin
aumentación, reflexión de canales y ruido en Fourier. Por tanto los efectos de subbanda,
de la ventana de \seg{3} y de la ausencia de aumentación no se evalúan en LaBraM ni en
EEGMamba.


